
\subsubsection{Training Corpus}

RedPajama-Data-V2 (English, snapshot 2023-14) serves as the source corpus, with pre-computed ccnet\_perplexity quality signals.

\subsubsection{Filtering Strategies}

Three strategies are applied to the same source:
\begin{itemize}
\item \textbf{Perplexity-filtered}: Bottom 30\% by ccnet\_perplexity
\item \textbf{Random-sampled}: Uniform random selection
\item \textbf{Inverse-perplexity}: Top 30\% by perplexity (control)
\end{itemize}

\subsubsection{Evaluation Benchmarks}

\begin{table}[htbp]
\centering
\resizebox{\columnwidth}{!}{%
\begin{tabular}{ll}
\toprule
Benchmark & Purpose \\
\midrule
MMLU & Primary contamination target \\
MMLU-Redux (2024+) & Time-stratified clean control \\
\bottomrule
\end{tabular}
}
\end{table}

\subsubsection{Model Configuration}

\begin{table}[htbp]
\centering
\resizebox{\columnwidth}{!}{%
\begin{tabular}{ll}
\toprule
Parameter & Value \\
\midrule
Model & Pythia-1B (planned); Pythia-70m (PoC) \\
Optimizer & AdamW \\
Learning rate & 2.5 $\times 10^{-4}$ (cosine decay) \\
Warmup & 1\% of steps \\
Batch size & 512 $\times$ 2048 tokens \\
Weight decay & 0.1 \\
\bottomrule
\end{tabular}
}
\end{table}

\subsubsection{Execution Mode}

All experiments were executed in methodology validation mode due to GPU infrastructure constraints (CUDA driver incompatibility). Results validate the measurement methodology using simulated contamination rather than naturally occurring perplexity-filtering differences. Full empirical validation with model training is pending GPU cluster access.
